% Congruences from self-conjugacy: Z8 (C5), Z9 (C13), Z10 (C14), Z22 (C23).
% Open questions for the author:
% - Credit for Z8 adds Arasu-Seberry 1998, Thm 3 (the lambda = 0 analogue, read by the refuter
%   in R4) and Arasu-Gordon-Zhang Thm 2.1 (Lander's form, in solutions.md's Origin); both added
%   to refs.bib. AGZ is cited as its arXiv preprint; add the journal version if wanted.
% - The worked examples (N173, N326, N514, N1747) were recomputed by hand and by a script from
%   the certificates in the Lean row files; they are not generated by numbers.py.

\section{Congruences from self-conjugacy}\label{sec:congruences}

The four criteria of this section share one setting. Let $\phi \colon G \to H$ be a surjective
homomorphism onto a nontrivial group $H$ of order $N$, with fibres of size $c = v/N$, and let
$p$ be a prime with
\begin{equation}\label{eq:sc}
  p \nmid N, \qquad p^i \equiv -1 \pmod{\exp H}, \qquad p^{2j} \mid k - \lambda
\end{equation}
for integers $i \ge 1$ and $j \ge 0$. The first two conditions say that $p$ is self-conjugate modulo
$\exp H$ (Definition~\ref{def:selfconj}) and prime to $|H|$; the second is the same as
$p^i h = -h$ for every $h \in H$, which is how Lean states it. Under~\eqref{eq:sc} the coset
sums of an SDS agree modulo $p^j$ (Lemma~\ref{lem:sc-congr}). The criteria then show that no
integer vector in one residue class modulo $p^j$ (or modulo a product of such prime powers)
has the bound, sum and square sum of Lemma~\ref{lem:contraction}. Each criterion assumes
$k \neq \lambda$.

\subsection{The congruence}

\begin{lemma}\label{lem:sc-congr}
Let $H \neq 0$ be a finite abelian group, $p$ a prime with $p \nmid |H|$ and $p^i h = -h$ for
all $h \in H$, for some $i \ge 1$. Let $b \colon H \to \Z$ and $n, \mu \in \Z$ with
\[
  \sum_{h} b_h\, b_{h+g} = n\,[g = 0] + \mu \qquad \text{for every } g \in H.
\]
If $p^{2j} \mid n$, then $b_h \equiv b_{h'} \pmod{p^j}$ for all $h, h' \in H$.
\end{lemma}

\begin{proof}
Induction on $j$; the case $j = 0$ is empty. Let $j \ge 1$.

\emph{Step 1: $b$ is constant modulo $p$.} Reduce modulo $p$ and write $B_h \in \mathbb{F}_p$
for $b_h$, $S = \sum_h B_h$ and $t = S/|H|$, which exists since $p \nmid |H|$. Since $p \mid n$,
all autocorrelations of $B$ equal $\mu$, and summing them over $g$ gives $S^2 = |H| \mu$, so
$\mu = |H| t^2$. Hence $B' = B - t$ has
$\sum_h B'_h B'_{h+g} = \mu - 2tS + |H| t^2 = 0$ for every $g$. In the group algebra
$\mathbb{F}_p[H]$ put $y = \sum_h B'_h\, [h]$ and $\bar y = \sum_h B'_h\, [-h]$. Then
$y \bar y = \sum_g \bigl(\sum_h B'_h B'_{h+g}\bigr) [g] = 0$. In characteristic $p$ the
Frobenius map gives $y^{p^i} = \sum_h B'_h\, [p^i h] = \bar y$, hence $y^{1+p^i} = 0$ and
$y^{p^{2i}} = \bar y^{\,p^i} = y$. As $p^{2i} \ge 1 + p^i$,
$y = y^{p^{2i}} = y^{1 + p^i} y^{p^{2i} - 1 - p^i} = 0$, so $B' = 0$.

\emph{Step 2: induction.} By Step~1, $b_h = t + p\, b'_h$ with $t = b_0$ and $b'$ integral.
Expanding, with $S = \sum_h b_h$,
\[
  p^2 \sum_h b'_h\, b'_{h+g} = \sum_h (b_h - t)(b_{h+g} - t) = n\,[g = 0] + \mu - 2tS + |H| t^2 .
\]
Evaluating at some $g_0 \neq 0$ shows that $\mu' = (\mu - 2tS + |H|t^2)/p^2$ is an integer, so
$b'$ has autocorrelations $(n/p^2)[g = 0] + \mu'$ with $p^{2(j-1)} \mid n/p^2$. By induction
$b'_h \equiv b'_{h'} \pmod{p^{j-1}}$, so $b_h \equiv b_{h'} \pmod{p^j}$.
\end{proof}

The coset sums of an SDS satisfy the hypothesis with $n = k - \lambda$ and $\mu = \lambda c$
(Lemma~\ref{lem:contraction}, (3) and~(4)). Hence, under~\eqref{eq:sc}, $a_h \equiv a_{h'}
\pmod{p^j}$ for all $h, h' \in H$ (\lean{SignedDifferenceSets.IsSDS.contract_congr_pow}; the
lemma itself is \lean{SignedDifferenceSets.congr_pow_of_autocorr}).

\begin{remark}
Lemma~\ref{lem:sc-congr} is Turyn's lemma in another form. For a nontrivial character $\chi$ of
$H$ the hypothesis gives $|\chi(b)|^2 = n$, so $p^j \mid \chi(b)$ by Lemma~\ref{lem:turyn}
with $e = \exp H$. Fourier inversion gives
$|H|\, b_h - \sum_{h'} b_{h'} = \sum_{\chi \neq \chi_0} \chi(b)\, \overline{\chi(h)}$, which
lies in $p^j \Z[\zeta_e] \cap \Z = p^j \Z$, and $|H|$ is a unit modulo $p$. The proof above
avoids algebraic number theory; it is the one checked in Lean.
\end{remark}

\subsection{Self-conjugacy on a quotient}

\begin{theorem}\label{thm:z8}
Let $G$ be a finite abelian group of order $v$, $k \ge 0$ and $\lambda$ integers with
$k \neq \lambda$, and $\phi \colon G \to H$ surjective onto a nontrivial group $H$, with
$v = c\,|H|$. Let $p$ be a prime and $i \ge 1$, $j \ge 0$ integers satisfying~\eqref{eq:sc}.
If $2c < p^j$, then $G$ has no SDS with parameters $(v,k,\lambda)$.
\end{theorem}

\begin{proof}
Let $D$ be such an SDS. By Lemma~\ref{lem:sc-congr} its coset sums agree modulo $p^j$, and
$|a_h - a_{h'}| \le 2c < p^j$ by Lemma~\ref{lem:contraction}(1), so all $a_h$ equal one integer
$a$. For $g \neq 0$ in $H$, Lemma~\ref{lem:contraction} gives
$\lambda c = \sum_h a_h a_{h+g} = |H|\, a^2 = \sum_h a_h^2 = k + \lambda(c - 1)$, so
$k = \lambda$, a contradiction.
\end{proof}

The case $H = G$, $c = 1$ is the plain self-conjugacy test on $D$ itself. The inequality
cannot be relaxed to $2c \le p^j$: for odd $v$, the map $D$ with $D(0) = -1$ and $D(x) = 1$
otherwise is a $(v, v, v - 4)$ SDS, and if $2$ is self-conjugate modulo $\exp G$, then $H = G$,
$p = 2$, $j = 1$ satisfy~\eqref{eq:sc} with $2c = p^j$.

\begin{example}
The case N173, $(36, 29, 4)$ in $C_3 \times C_{12}$, has $n = 25$. Take $H = C_3 \times C_6$
($|H| = 18$, $c = 2$) and $p = 5$, $i = j = 1$: $5 \equiv -1 \pmod 6$, $5 \nmid 18$ and
$25 \mid n$. The coset sums lie in $[-2, 2]$ and agree modulo $5$, so they are all equal, and
Theorem~\ref{thm:z8} applies since $2c = 4 < 5$.
\end{example}

\begin{remark}
Theorem~\ref{thm:z8} applies Turyn's self-conjugacy lemma~\cite[Lemma~3]{turyn}, in Lander's
group-ring form~\cite[Thm.~2.1]{arasu-gordon-zhang}, to the coset sums of an SDS; this form
with the bound $2c < p^j$ is ours. He and Wu~\cite[Thm.~3.2]{he-wu} use self-conjugacy only
for the odd-exponent case: if $p^{i} \equiv -1 \pmod m$ for some $m > 1$ dividing $\exp G$,
then the exponent of $p$ in $n$ is even. For circulant weighing matrices ($\lambda = 0$) the
corresponding bound, with $p^j \le c$ in place of $p^j \le 2c$, is~\cite[Thm.~3]{arasu-seberry},
quoted there from earlier work of the same authors. The criterion settles \zrows{8} cases. Lean:
\lean{SignedDifferenceSets.not_hasSDS_of_selfConj}; the cases are checked in the modules
\lean{SignedDifferenceSets.SelfConjugacyRows.Part01} to \lean{Part13}.
\end{remark}

\subsection{Moments of the coset sums}
When $2c \ge p^j$ the congruence leaves several values per residue class. Let
$q > 0$ be the modulus, $N = |H|$, $Q = k + \lambda(c - 1)$, and suppose all coset sums are
congruent to $t_0 \in \{0, \dots, q - 1\}$. Writing $a_h = t_0 + q\, y_h$ and $s' = \sum_h a_h$,
the integers $y_h$ have
\begin{equation}\label{eq:moments}
  S_1 = \sum_h y_h = \frac{s' - N t_0}{q}, \qquad
  S_2 = \sum_h y_h^2 = \frac{Q - 2 t_0 s' + N t_0^2}{q^2}.
\end{equation}

\begin{definition}\label{def:moments}
Let $N, q > 0$ and $s', Q, t_0$ be integers. The tuple $(N, q, s', Q, t_0)$ \emph{passes the
moments test} if $q \mid s' - N t_0$ and $q^2 \mid Q - 2 t_0 s' + N t_0^2$, and, with $S_1$,
$S_2$ as in~\eqref{eq:moments}, $m = \lfloor S_1 / N \rfloor$, $A = S_1 - N m$ and
$B = S_2 - 2 m S_1 + N m^2$,
\[
  S_2 \equiv S_1 \pmod 2, \qquad A \le B, \qquad B \neq 0 .
\]
\end{definition}

\begin{lemma}\label{lem:moments}
Let $a \colon H \to \Z$ take values congruent to $t_0$ modulo $q > 0$, not all equal. Then
$(|H|, q, \sum_h a_h, \sum_h a_h^2, t_0)$ passes the moments test.
\end{lemma}

\begin{proof}
With $a_h = t_0 + q y_h$ the divisibilities hold and~\eqref{eq:moments} gives $S_1, S_2$. Put
$z_h = y_h - m$; then $\sum_h z_h = A$ and $\sum_h z_h^2 = B$. Each $z_h^2 - z_h = z_h(z_h - 1)$
is even and nonnegative, so $A \le B$ and $B - A$ is even; as $B - A \equiv S_2 - S_1 \pmod 2$
($m(m+1)$ is even), $S_2 \equiv S_1 \pmod 2$. If $B = 0$, all $y_h$ equal $m$ and all $a_h$
are equal.
\end{proof}

\begin{theorem}\label{thm:z9}
Let $G$, $k \neq \lambda$, $\phi \colon G \to H$, $c$ and $p, i, j$ be as in
Theorem~\ref{thm:z8}, without the condition $2c < p^j$, and let $s$ be an integer with
$s^2 = k + \lambda(v - 1)$. If for both $s' \in \{s, -s\}$ and every
$t_0 \in \{0, \dots, p^j - 1\}$ the tuple $(|H|, p^j, s', k + \lambda(c-1), t_0)$ fails the
moments test, then $G$ has no SDS with parameters $(v,k,\lambda)$.
\end{theorem}

\begin{proof}
Let $D$ be such an SDS. By Lemma~\ref{lem:sc-congr} all coset sums are congruent modulo $p^j$
to the residue $t_0 \in \{0, \dots, p^j - 1\}$ of $a_0$. They are not all equal, as the proof of
Theorem~\ref{thm:z8} shows from $k \neq \lambda$. By Lemma~\ref{lem:contraction},
$\sum_h a_h = \pm s$ and $\sum_h a_h^2 = k + \lambda(c - 1)$, so Lemma~\ref{lem:moments}
contradicts the hypothesis.
\end{proof}

\begin{example}
The case N326, $(48, 27, 2)$ in $C_2 \times C_{24}$, has $n = 25$ and $s = 11$. Take
$H = C_2 \times C_6$ ($N = 12$, $c = 4$) and $p = 5$, $i = j = 1$; here $2c = 8 > 5$, so
Theorem~\ref{thm:z8} does not apply, and $Q = 33$. For $s' = 11$ the divisibility
$5 \mid 11 - 12 t_0$ forces $t_0 = 3$, and then $S_1 = -5$, $S_2 = 3$. For $s' = -11$ it forces
$t_0 = 2$, and $S_1 = -7$, $S_2 = 5$. Both fail, since $\sum_h y_h^2 \ge \sum_h |y_h| \ge
|S_1|$ for integers; in the terms of Definition~\ref{def:moments}, $m = -1$ and $B = 5 < 7 = A$,
respectively $B = 3 < 5 = A$.
\end{example}

\begin{remark}
The identities behind~\eqref{eq:moments} are Gordon's intersection numbers~\cite[Lemma~5.2]{gordon}
(also~\cite[Lemma~4.1]{he-wu}); the test built on them is ours. The criterion settles
\zrows{9} further cases. Lean: \lean{SignedDifferenceSets.not_hasSDS_of_moments}; the cases are
checked in the same modules as those of Theorem~\ref{thm:z8}.
\end{remark}

\subsection{The box bound}
The moments test ignores the bound $|a_h| \le c$. With it, $q y_h = a_h - t_0$ lies in
$[-c - t_0, c - t_0]$, so $\mathrm{lo} \le y_h \le \mathrm{hi}$ with
\[
  \mathrm{lo} = -\left\lfloor \frac{c + t_0}{q} \right\rfloor, \qquad
  \mathrm{hi} = \left\lfloor \frac{c - t_0}{q} \right\rfloor .
\]

\begin{definition}\label{def:box}
Let $c \ge 0$ be an integer. The tuple $(N, q, s', Q, t_0)$ \emph{passes the box test with $c$}
if it passes the moments test and, with $S_1, S_2$ as in~\eqref{eq:moments},
\[
  S_2 \le (\mathrm{lo} + \mathrm{hi})\, S_1 - N\, \mathrm{lo}\, \mathrm{hi}.
\]
\end{definition}

\begin{lemma}\label{lem:box}
If $a$ is as in Lemma~\ref{lem:moments} and moreover $|a_h| \le c$ for all $h$, then
$(|H|, q, \sum_h a_h, \sum_h a_h^2, t_0)$ passes the box test with $c$.
\end{lemma}

\begin{proof}
Each term of $\sum_h (y_h - \mathrm{lo})(\mathrm{hi} - y_h)$ is nonnegative, and the sum equals
$(\mathrm{lo} + \mathrm{hi}) S_1 - S_2 - N\,\mathrm{lo}\,\mathrm{hi}$.
\end{proof}

\begin{theorem}\label{thm:z10}
Let $G$, $k \neq \lambda$, $\phi$, $c$, $p, i, j$ and $s$ be as in Theorem~\ref{thm:z9}. If
for both $s' \in \{s, -s\}$ and every $t_0 \in \{0, \dots, p^j - 1\}$ the tuple
$(|H|, p^j, s', k + \lambda(c-1), t_0)$ fails the box test with $c$, then $G$ has no SDS with
parameters $(v,k,\lambda)$.
\end{theorem}

\begin{proof}
As for Theorem~\ref{thm:z9}, with Lemma~\ref{lem:box} in place of Lemma~\ref{lem:moments};
the bound $|a_h| \le c$ is Lemma~\ref{lem:contraction}(1).
\end{proof}

\begin{example}
The case N514, $(56, 44, 8)$ in $C_{56}$, has $n = 36$ and $s = 22$. Take $H = C_{28}$
($N = 28$, $c = 2$, $Q = 52$), $p = 3$, $j = 1$, $i = 3$ ($27 \equiv -1 \pmod{28}$). For
$s' = 22$ only $t_0 = 1$ survives the divisibilities, with $S_1 = -2$, $S_2 = 4$; the moments
test passes ($m = -1$, $A = 26 \le B = 28$). But $\mathrm{lo} = -1$, $\mathrm{hi} = 0$, so
$y_h^2 = -y_h$ and $S_2 = -S_1 = 2$ would be needed: the bound reads $4 \le 2$. In other words,
coset sums in $\{1, -2\}$ with sum $22$ have square sum $34 \ne 52$. For $s' = -22$ only
$t_0 = 2$ survives, with $S_1 = -26$, $S_2 = 28$, and the bound reads $28 \le 26$.
\end{example}

\begin{remark}
The box bound is ours. It settles \zrows{10} further cases. Lean:
\lean{SignedDifferenceSets.not_hasSDS_of_box}; the cases are checked in the modules
\lean{SignedDifferenceSets.BoxRows.Part01} to \lean{Part03}.
\end{remark}

\subsection{A composite modulus}
Several primes can satisfy~\eqref{eq:sc} on the same quotient. Their moduli are coprime, so
the congruences combine.

\begin{theorem}\label{thm:z22}
Let $G$, $k \neq \lambda$, $\phi \colon G \to H$, $c$ and $s$ be as in Theorem~\ref{thm:z9}.
Let $P$ be a finite set of primes and, for each $p \in P$, let $i_p \ge 1$ and $j_p \ge 0$ be
integers such that $p, i_p, j_p$ satisfy~\eqref{eq:sc}. Put $q = \prod_{p \in P} p^{j_p}$. If
for both $s' \in \{s, -s\}$ and every $t_0 \in \{0, \dots, q - 1\}$ the tuple
$(|H|, q, s', k + \lambda(c-1), t_0)$ fails the box test with $c$, then $G$ has no SDS with
parameters $(v,k,\lambda)$.
\end{theorem}

\begin{proof}
For each $p \in P$, Lemma~\ref{lem:sc-congr} gives $p^{j_p} \mid a_h - a_{h'}$. The moduli
are pairwise coprime, so $q \mid a_h - a_{h'}$. Lemmas~\ref{lem:moments} and~\ref{lem:box}
hold for any modulus $q > 0$, and the proof of Theorem~\ref{thm:z10} goes through.
\end{proof}

\begin{example}
The case N1747, $(95, 56, 20)$ in $C_{95}$, has $n = 36$ and $s = 44$. On $H = C_{19}$
($c = 5$, $Q = 136$) both $p = 2$ and $p = 3$ satisfy~\eqref{eq:sc} with $i = 9$ and $j = 1$
($2^9 \equiv 3^9 \equiv -1 \pmod{19}$). Alone, each passes the box test for some $t_0$: $q = 2$
with $t_0 = 0$, $q = 3$ with $t_0 = 2$ ($s' = 44$) or $t_0 = 1$ ($s' = -44$). Together,
$q = 6$. For $s' = 44$ only $t_0 = 2$ survives the divisibilities; the coset sums would lie in
$\{2, -4\}$, so their sum is at most $38 < 44$ (in the test: $S_1 = 1$ but $\mathrm{hi} = 0$).
For $s' = -44$ only $t_0 = 4$ survives; the coset sums would lie in $\{4, -2\}$, and the sum
$-44$ needs $20$ of the $19$ coset sums to be $-2$ (in the test: $S_2 = 22 > 20 = -S_1$).
\end{example}

\begin{remark}
Combining the primes is ours, on Turyn's self-conjugacy. The criterion settles \zrows{22}
further cases. Lean: \lean{SignedDifferenceSets.not_hasSDS_of_compositeModulus}; the cases are
checked in the module \lean{SignedDifferenceSets.CompositeModulusRows}.
\end{remark}
